\documentclass[10pt,a4paper]{amsart}
\usepackage[margin=19mm]{geometry}
\usepackage{amsmath,amssymb,mathtools}
\usepackage[scr=rsfs,cal=euler]{mathalfa}
\usepackage{xfrac}
\usepackage[unicode,hidelinks,bookmarks=false]{hyperref}
\newcommand{\ie}{i.e.\ }
\newcommand{\cf}{cf.\ }
\newcommand{\resp}{respectively\ }
\newcommand{\Subsection}{\subsection}
\providecommand{\nobreakdash}{\nobreak\hbox{-}}
\setlength{\emergencystretch}{2em}
\multlinegap0pt
\usepackage{enumerate}
\usepackage{tikz}
\newtheorem{lemma}{Lemma}
\title{Can we further improve the central limit theorem for stationary $\rho$-mixing sequences $?$}
\author{Muneya Matsui}
\date{Source: arXiv:2609.04638v1 (CC BY 4.0)}
\begin{document}
\maketitle

\noindent\emph{Source and licence.} Can we further improve the central limit theorem for stationary $\rho$-mixing sequences $?$. Version arXiv:2609.04638v1 (CC BY 4.0), distributed by arXiv under the Creative Commons Attribution 4.0 International licence, http://creativecommons.org/licenses/by/4.0/, as recorded in the arXiv licence metadata for this exact version and shown on the abstract page https://arxiv.org/abs/2609.04638v1. Full licence notice: \texttt{licenses/arxiv-2609.04638-CC-BY-4.0.txt}.\par\smallskip

\noindent\textit{Editorial scope.} Counted unit: the article's lemma aux:lem:ineq (source0.tex lines 1427--1434). The article's complete original proof of it is delivered with it (source0.tex lines 1436--1478, one \texttt{proof} environment of 43 lines). No mathematics is written, repaired or re-derived: the statement and the proof are the article's own lines, in the article's own notation, and no retained line is substituted or re-flowed.  No further source-local context is needed: every cross-reference of every retained line resolves inside the two delivered spans.  Nothing was omitted from any retained span, so the package carries no omission record.  No retained line contains a citation, so the wrapper carries no bibliography and no bibliography entry.  No retained line contains a figure environment, an image-inclusion command or drawing code, so no image is delivered and the package declares no asset dependency.  Editorial formatting only, in addition: an AMS \texttt{amsart} preamble that restates the article's own declarations and macros used by the retained lines, namely the environments \texttt{lemma} on separate counters, together with the standard packages those lines need.  The retained environments print in this document as Lemma 1 in order of appearance, whereas the article prints the same items with the numbering of its own full document.  The complete native source of the article as distributed in the arXiv e-print for this exact version is bundled unchanged as \texttt{sources/209378/source0.tex}.
\begin{lemma}
\label{aux:lem:ineq}
For all $x,y\ge 0$ and $a\ge 1$, the inequality
\begin{align*}
(x+y)^a \le x^a +y^a +a (2^{a-2}\vee 1) (x^{a-1}y +xy^{a-1})
\end{align*}
holds. 
\end{lemma}

\begin{proof}
The inequality trivially holds for $x=0$, so without loss of generality we may assume $x>0$. 
Dividing both sides by $x^a$ and replacing $y/x$ with $x$, it suffices to show that
\[
(1+x)^a \le 1+x^a + a(2^{a-2}\vee 1) (x+x^{a-1}),\quad x\ge 0.
\]
We define 
\[
f(x) = 1+ x^a +a(2^{a-2}\vee 1)(x+x^{a-1}) -(1+x)^a
\]
and show that $f(x)\ge 0$ for all $x\ge 0$. Observing the symmetry property 
$f(1/x)=x^{-a}f(x)$, it is enough to consider $f$ on either $x\in [0,1]$ or $x\in [1,\infty)$. 

First, consider the case $a\in (1,3]$ and $x\in [0,1]$. Note that $f(0)=0$ and $f(1)=2+2a(2^{a-2}\vee 1)-2^a>0$. The case $a=1$ is trivial. 
We analyze the sign of the derivative
\[
f'(x) = a x^{a-1} + a(2^{a-2}\vee 1)\{1+(a-1) x^{a-2}\}-a(1+x)^{a-1}. 
\]
When $a \in (1,2]$, the triangle inequality $(1+x)^{a-1} \le 1+ x^{a-1}$ immediately yields $f'(x) \ge 0$. 
For $a \in (2,3]$, by the convexity of $(1+x)^{a-1}$, we have 
\[
(1+x)^{a-1} \le 1 +(a-1)(1+x)^{a-2} x \le 1+(a-1) 2^{a-2} x,
\]
which implies that
\[
f'(x) \ge a x^{a-1}+ a(2^{a-2}-1)+a(a-1) 2^{a-2} (x^{a-2}-x) >0.
\]
Hence, $f(x)\ge 0$ is proved for $a\in (1,3]$.

Next, suppose $a \ge 3$. We examine $f$ on $x\in [1,\infty)$, and observe that 
$f(1)= 2+a2^{a-1}-2^a>0$ and $f(\infty)=\infty$ (by applying L'H\^opital's rule to $f(1/x)=x^{-a}f(x)$). 
Moreover, $f'(x)\ge 0$ holds on $[1,\infty)$. Indeed, applying the mean value theorem to $(1+z)^{a-1}$ for $z\in [0,1]$, we obtain 
\begin{align*}
f'(x) &= a x^{a-1} \Big\{
1+2^{a-2}\big(x^{1-a}+(a-1)x^{-1}\big)-(1+x^{-1})^{a-1}
\Big\} \\
&\ge a x^{a-1} \Big\{
(a-1) \big(2^{a-2} -(1+x_\ast^{-1})^{a-2}\big)x^{-1} + 2^{a-2}x^{1-a}
\Big\} \\
& \ge a 2^{a-2}, 
\end{align*} 
where $x_\ast$ is some value in $[1,\infty)$. This completes the proof.
\end{proof}

\end{document}
